% jcr-header.tex (jcrquarto 1.0.0)
% Preamble for JCR manuscripts. pdflatex-compatible fonts: Times for the text,
% tables, captions, and notes; Helvetica is loaded for any sans serif text.

% ---- Fonts ----
\usepackage{mathptmx}
\usepackage[scaled=0.92]{helvet}

% ---- Quarto 1.4 (pandoc before 3.2) ----
% Those versions wrap \paragraph and \subparagraph in a way that titlesec
% cannot read ("Argument of \paragraph has an extra }"). The original
% commands are put back before titlesec is loaded; the heading formats below
% define them anew.
\makeatletter
\@ifundefined{xxxParagraphStar}{%
  \@ifundefined{oldparagraph}{}{\let\paragraph\oldparagraph}%
  \@ifundefined{oldsubparagraph}{}{\let\subparagraph\oldsubparagraph}%
}{}
\makeatother

% ---- Packages ----
\usepackage{float}
\usepackage{setspace}
\usepackage{ragged2e}
\usepackage{titlesec}
\usepackage{etoolbox}
\usepackage{longtable}
\usepackage{booktabs}
\usepackage[bottom]{footmisc}
\usepackage{indentfirst}
\usepackage{fancyhdr}

% ---- Characters typed directly ----
% pdfLaTeX has no glyphs of its own for Greek letters and for many signs that
% are typed as characters (a Greek alpha, a minus sign, "less than or equal").
% In the text, Greek letters are taken upright from the Symbol font, as they
% appear in a Word file; in formulas they are the usual math letters. The
% other signs are set in math mode. XeLaTeX and LuaLaTeX read all of these
% characters without help, so the definitions apply to pdfLaTeX only.
\usepackage{iftex}
\usepackage{pifont}
\ifPDFTeX
\providecommand{\texorpdfstring}[2]{#1}
\newcommand{\jcrgreek}[3]{\texorpdfstring{\TextOrMath{\Pisymbol{psy}{#1}}{#2}}{#3}}
\newcommand{\jcrsign}[1]{\texorpdfstring{\ensuremath{#1}}{}}
\DeclareUnicodeCharacter{03B1}{\jcrgreek{97}{\alpha}{alpha}}
\DeclareUnicodeCharacter{03B2}{\jcrgreek{98}{\beta}{beta}}
\DeclareUnicodeCharacter{03B3}{\jcrgreek{103}{\gamma}{gamma}}
\DeclareUnicodeCharacter{03B4}{\jcrgreek{100}{\delta}{delta}}
\DeclareUnicodeCharacter{03B5}{\jcrgreek{101}{\varepsilon}{epsilon}}
\DeclareUnicodeCharacter{03F5}{\jcrgreek{101}{\epsilon}{epsilon}}
\DeclareUnicodeCharacter{03B6}{\jcrgreek{122}{\zeta}{zeta}}
\DeclareUnicodeCharacter{03B7}{\jcrgreek{104}{\eta}{eta}}
\DeclareUnicodeCharacter{03B8}{\jcrgreek{113}{\theta}{theta}}
\DeclareUnicodeCharacter{03D1}{\jcrgreek{74}{\vartheta}{theta}}
\DeclareUnicodeCharacter{03B9}{\jcrgreek{105}{\iota}{iota}}
\DeclareUnicodeCharacter{03BA}{\jcrgreek{107}{\kappa}{kappa}}
\DeclareUnicodeCharacter{03BB}{\jcrgreek{108}{\lambda}{lambda}}
\DeclareUnicodeCharacter{03BC}{\jcrgreek{109}{\mu}{mu}}
\DeclareUnicodeCharacter{03BD}{\jcrgreek{110}{\nu}{nu}}
\DeclareUnicodeCharacter{03BE}{\jcrgreek{120}{\xi}{xi}}
\DeclareUnicodeCharacter{03BF}{\jcrgreek{111}{o}{o}}
\DeclareUnicodeCharacter{03C0}{\jcrgreek{112}{\pi}{pi}}
\DeclareUnicodeCharacter{03D6}{\jcrgreek{118}{\varpi}{pi}}
\DeclareUnicodeCharacter{03C1}{\jcrgreek{114}{\rho}{rho}}
\DeclareUnicodeCharacter{03C2}{\jcrgreek{86}{\varsigma}{sigma}}
\DeclareUnicodeCharacter{03C3}{\jcrgreek{115}{\sigma}{sigma}}
\DeclareUnicodeCharacter{03C4}{\jcrgreek{116}{\tau}{tau}}
\DeclareUnicodeCharacter{03C5}{\jcrgreek{117}{\upsilon}{upsilon}}
\DeclareUnicodeCharacter{03C6}{\jcrgreek{106}{\varphi}{phi}}
\DeclareUnicodeCharacter{03D5}{\jcrgreek{102}{\phi}{phi}}
\DeclareUnicodeCharacter{03C7}{\jcrgreek{99}{\chi}{chi}}
\DeclareUnicodeCharacter{03C8}{\jcrgreek{121}{\psi}{psi}}
\DeclareUnicodeCharacter{03C9}{\jcrgreek{119}{\omega}{omega}}
\DeclareUnicodeCharacter{0391}{\jcrgreek{65}{\mathrm{A}}{A}}
\DeclareUnicodeCharacter{0392}{\jcrgreek{66}{\mathrm{B}}{B}}
\DeclareUnicodeCharacter{0393}{\jcrgreek{71}{\Gamma}{Gamma}}
\DeclareUnicodeCharacter{0394}{\jcrgreek{68}{\Delta}{Delta}}
\DeclareUnicodeCharacter{2206}{\jcrgreek{68}{\Delta}{Delta}}
\DeclareUnicodeCharacter{0395}{\jcrgreek{69}{\mathrm{E}}{E}}
\DeclareUnicodeCharacter{0396}{\jcrgreek{90}{\mathrm{Z}}{Z}}
\DeclareUnicodeCharacter{0397}{\jcrgreek{72}{\mathrm{H}}{H}}
\DeclareUnicodeCharacter{0398}{\jcrgreek{81}{\Theta}{Theta}}
\DeclareUnicodeCharacter{0399}{\jcrgreek{73}{\mathrm{I}}{I}}
\DeclareUnicodeCharacter{039A}{\jcrgreek{75}{\mathrm{K}}{K}}
\DeclareUnicodeCharacter{039B}{\jcrgreek{76}{\Lambda}{Lambda}}
\DeclareUnicodeCharacter{039C}{\jcrgreek{77}{\mathrm{M}}{M}}
\DeclareUnicodeCharacter{039D}{\jcrgreek{78}{\mathrm{N}}{N}}
\DeclareUnicodeCharacter{039E}{\jcrgreek{88}{\Xi}{Xi}}
\DeclareUnicodeCharacter{039F}{\jcrgreek{79}{\mathrm{O}}{O}}
\DeclareUnicodeCharacter{03A0}{\jcrgreek{80}{\Pi}{Pi}}
\DeclareUnicodeCharacter{03A1}{\jcrgreek{82}{\mathrm{P}}{P}}
\DeclareUnicodeCharacter{03A3}{\jcrgreek{83}{\Sigma}{Sigma}}
\DeclareUnicodeCharacter{03A4}{\jcrgreek{84}{\mathrm{T}}{T}}
\DeclareUnicodeCharacter{03A5}{\jcrgreek{85}{\Upsilon}{Upsilon}}
\DeclareUnicodeCharacter{03A6}{\jcrgreek{70}{\Phi}{Phi}}
\DeclareUnicodeCharacter{03A7}{\jcrgreek{67}{\mathrm{X}}{X}}
\DeclareUnicodeCharacter{03A8}{\jcrgreek{89}{\Psi}{Psi}}
\DeclareUnicodeCharacter{03A9}{\jcrgreek{87}{\Omega}{Omega}}
\DeclareUnicodeCharacter{2212}{\texorpdfstring{\ensuremath{-}}{-}}
\DeclareUnicodeCharacter{2264}{\jcrsign{\leq}}
\DeclareUnicodeCharacter{2265}{\jcrsign{\geq}}
\DeclareUnicodeCharacter{2260}{\jcrsign{\neq}}
\DeclareUnicodeCharacter{2248}{\jcrsign{\approx}}
\DeclareUnicodeCharacter{2261}{\jcrsign{\equiv}}
\DeclareUnicodeCharacter{223C}{\jcrsign{\sim}}
\DeclareUnicodeCharacter{2245}{\jcrsign{\cong}}
\DeclareUnicodeCharacter{221D}{\jcrsign{\propto}}
\DeclareUnicodeCharacter{2213}{\jcrsign{\mp}}
\DeclareUnicodeCharacter{2217}{\jcrsign{\ast}}
\DeclareUnicodeCharacter{22C5}{\jcrsign{\cdot}}
\DeclareUnicodeCharacter{2219}{\jcrsign{\bullet}}
\DeclareUnicodeCharacter{221A}{\jcrsign{\surd}}
\DeclareUnicodeCharacter{221E}{\jcrsign{\infty}}
\DeclareUnicodeCharacter{2211}{\jcrsign{\sum}}
\DeclareUnicodeCharacter{220F}{\jcrsign{\prod}}
\DeclareUnicodeCharacter{222B}{\jcrsign{\int}}
\DeclareUnicodeCharacter{2202}{\jcrsign{\partial}}
\DeclareUnicodeCharacter{2207}{\jcrsign{\nabla}}
\DeclareUnicodeCharacter{2208}{\jcrsign{\in}}
\DeclareUnicodeCharacter{2209}{\jcrsign{\notin}}
\DeclareUnicodeCharacter{2282}{\jcrsign{\subset}}
\DeclareUnicodeCharacter{2286}{\jcrsign{\subseteq}}
\DeclareUnicodeCharacter{2229}{\jcrsign{\cap}}
\DeclareUnicodeCharacter{222A}{\jcrsign{\cup}}
\DeclareUnicodeCharacter{2205}{\jcrsign{\emptyset}}
\DeclareUnicodeCharacter{2200}{\jcrsign{\forall}}
\DeclareUnicodeCharacter{2203}{\jcrsign{\exists}}
\DeclareUnicodeCharacter{2227}{\jcrsign{\wedge}}
\DeclareUnicodeCharacter{2228}{\jcrsign{\vee}}
\DeclareUnicodeCharacter{2032}{\jcrsign{{}^{\prime}}}
\DeclareUnicodeCharacter{2033}{\jcrsign{{}^{\prime\prime}}}
\DeclareUnicodeCharacter{2194}{\jcrsign{\leftrightarrow}}
\DeclareUnicodeCharacter{21D2}{\jcrsign{\Rightarrow}}
\DeclareUnicodeCharacter{21D0}{\jcrsign{\Leftarrow}}
\DeclareUnicodeCharacter{21D4}{\jcrsign{\Leftrightarrow}}
\DeclareUnicodeCharacter{2197}{\jcrsign{\nearrow}}
\DeclareUnicodeCharacter{2198}{\jcrsign{\searrow}}
\DeclareUnicodeCharacter{25CF}{\jcrsign{\bullet}}
\DeclareUnicodeCharacter{25CB}{\jcrsign{\circ}}
\DeclareUnicodeCharacter{25A0}{\jcrsign{\blacksquare}}
\DeclareUnicodeCharacter{25A1}{\jcrsign{\square}}
\DeclareUnicodeCharacter{25B2}{\jcrsign{\blacktriangle}}
\DeclareUnicodeCharacter{25B3}{\jcrsign{\triangle}}
\DeclareUnicodeCharacter{25BC}{\jcrsign{\blacktriangledown}}
\DeclareUnicodeCharacter{2605}{\jcrsign{\bigstar}}
\DeclareUnicodeCharacter{2606}{\texorpdfstring{\ding{73}}{}}
\DeclareUnicodeCharacter{2023}{\jcrsign{\triangleright}}
\DeclareUnicodeCharacter{2713}{\texorpdfstring{\ding{51}}{}}
\DeclareUnicodeCharacter{2714}{\texorpdfstring{\ding{52}}{}}
\DeclareUnicodeCharacter{2717}{\texorpdfstring{\ding{55}}{}}
\DeclareUnicodeCharacter{2718}{\texorpdfstring{\ding{56}}{}}
\DeclareUnicodeCharacter{2070}{\textsuperscript{0}}
\DeclareUnicodeCharacter{2074}{\textsuperscript{4}}
\DeclareUnicodeCharacter{2075}{\textsuperscript{5}}
\DeclareUnicodeCharacter{2076}{\textsuperscript{6}}
\DeclareUnicodeCharacter{2077}{\textsuperscript{7}}
\DeclareUnicodeCharacter{2078}{\textsuperscript{8}}
\DeclareUnicodeCharacter{2079}{\textsuperscript{9}}
\DeclareUnicodeCharacter{207A}{\textsuperscript{+}}
\DeclareUnicodeCharacter{207B}{\textsuperscript{\ensuremath{-}}}
\DeclareUnicodeCharacter{2080}{\textsubscript{0}}
\DeclareUnicodeCharacter{2081}{\textsubscript{1}}
\DeclareUnicodeCharacter{2082}{\textsubscript{2}}
\DeclareUnicodeCharacter{2083}{\textsubscript{3}}
\DeclareUnicodeCharacter{2084}{\textsubscript{4}}
\DeclareUnicodeCharacter{2085}{\textsubscript{5}}
\DeclareUnicodeCharacter{2086}{\textsubscript{6}}
\DeclareUnicodeCharacter{2087}{\textsubscript{7}}
\DeclareUnicodeCharacter{2088}{\textsubscript{8}}
\DeclareUnicodeCharacter{2089}{\textsubscript{9}}
\DeclareUnicodeCharacter{2153}{\textsuperscript{1}/\textsubscript{3}}
\DeclareUnicodeCharacter{2154}{\textsuperscript{2}/\textsubscript{3}}
\DeclareUnicodeCharacter{2060}{\nobreak}
\DeclareUnicodeCharacter{FB00}{ff}
\DeclareUnicodeCharacter{FB01}{fi}
\DeclareUnicodeCharacter{FB02}{fl}
\DeclareUnicodeCharacter{FB03}{ffi}
\DeclareUnicodeCharacter{FB04}{ffl}
\fi

% ---- Tables: 12-point Times, as in JCR's sample manuscripts, single spaced ----
% \setstretch{1} rather than \singlespacing, which adds a blank line
\AtBeginEnvironment{table}{\normalsize\setstretch{1}}
\AtBeginEnvironment{table*}{\normalsize\setstretch{1}}
\AtBeginEnvironment{longtable}{\normalsize\setstretch{1}}
% A blank line above a table label, and less space between a table and its note
\setlength{\LTpre}{24pt}
\setlength{\LTpost}{4pt}
% A table that fits on a page is not split: \jcrkeep{height} starts a new
% page when less than that height is left on the current page
\newcommand{\jcrkeep}[1]{%
  \par\ifdim\pagetotal>0pt\ifdim\dimexpr\pagegoal-\pagetotal\relax<#1\relax\newpage\fi\fi}
% A table or a figure without a label takes no number. LaTeX counts every
% title, and pandoc 3.2 (in Quarto 1.5) counts every table, with or without
% one; the counter is put back afterward, so that the numbers agree with the
% Word file, and a title is printed without "TABLE 1" before it.
\newcounter{jcrsaved}
\newcounter{jcrplain}
\newcommand{\jcrplainbegin}[1]{\setcounter{jcrsaved}{\value{#1}}\stepcounter{jcrplain}%
  \begingroup\ifdefined\captionsetup\captionsetup{format=plain,labelformat=empty,labelsep=none}\fi%
  \expandafter\def\csname theH#1\endcsname{plain.\arabic{jcrplain}}}
\newcommand{\jcrplainend}[1]{\endgroup\setcounter{#1}{\value{jcrsaved}}}

% ---- Captions: "TABLE 1" or "FIGURE 1" centered, title on the next line ----
% The title is set in capitals by the jcrquarto filters, so captions can be
% typed in ordinary case in the .qmd file. Captions use the text font (Times)
% and the line spacing of the text, as in JCR's sample manuscripts.
\usepackage[labelsep=newline,labelfont={bf},textfont={},font={rm,normalsize,stretch=2},
            justification=centering,singlelinecheck=false,skip=4pt]{caption}
\captionsetup[table]{position=above}
\captionsetup[figure]{position=above}
\makeatletter
\renewcommand{\fnum@table}{\MakeUppercase{\tablename}~\thetable}
\renewcommand{\fnum@figure}{\MakeUppercase{\figurename}~\thefigure}
\makeatother
% The captions of the panels of a figure: "(a) Text", centered, in the type
% of the text. (Quarto loads the subcaption package when a document has a
% figure with panels.)
\makeatletter
\AtBeginDocument{\@ifpackageloaded{subcaption}{%
  \captionsetup[sub]{font={rm,normalsize,singlespacing},labelfont={},textfont={},
    labelformat=parens,labelsep=space,justification=centering,singlelinecheck=false,skip=4pt}}{}}
\makeatother
% The caption of a picture in a panel without a number: the same, without a letter
\newcommand{\jcrpanelcaption}[1]{{\captionsetup{font={rm,normalsize,singlespacing},labelfont={},
    textfont={},justification=centering,singlelinecheck=false,skip=4pt}\caption*{#1}}}

% ---- Equations: "put bold number flush right in parentheses" ----
\makeatletter
\AtBeginDocument{\@ifpackageloaded{amsmath}{%
  \def\tagform@#1{\maketag@@@{\textbf{(\ignorespaces#1\unskip\@@italiccorr)}}}%
}{%
  \renewcommand{\@eqnnum}{{\normalfont\normalcolor\textbf{(\theequation)}}}%
}}
\makeatother

% ---- Line break inside a heading (appendix label, then title) ----
% Protected, so that the capitalizing of primary headings leaves it alone
\protected\def\jcrbr{\\}

% ---- No underlining ----
\renewcommand{\underline}[1]{#1}

% ---- Table and figure notes: at the bottom, paragraph indented ----
% Written in the .qmd file as a ::: {.jcr-note} div; the filter adds "NOTE."
% 12-point type with a half-inch paragraph indent, as in JCR's sample
% manuscripts; single spaced, like the table.
\newenvironment{jcrnotes}{%
  \par\nopagebreak[4]\vskip 4pt%
  \begingroup\normalsize\setstretch{1}\raggedright
  \setlength{\parindent}{0.5in}%
}{%
  \par\endgroup\vskip 12pt%
}
% (jcrfigurenote: the name used by jcrquarto 1.0 for the same layout)
\newenvironment{jcrfigurenote}{\begin{jcrnotes}}{\end{jcrnotes}}

% ---- Hypotheses: "bold, indented, syllabus style" ----
% Written in the .qmd file as a ::: {.hypothesis} div. As in the style
% guide's example, the label ("H1:") is bold and stands at the paragraph
% indent, and the hypothesis is set as a block one inch from the margin.
% \addvspace keeps the space single between two hypotheses in a row.
\newlength{\jcrhypwidth}
\newenvironment{jcrhypothesis}{%
  \par\addvspace{12pt}\begingroup\setlength{\parindent}{0pt}%
  \setlength{\leftskip}{1in}%
}{%
  \par\endgroup\addvspace{12pt}%
}
% The label hangs half an inch to the left of the hypothesis
\newcommand{\jcrhyplabel}[1]{%
  \settowidth{\jcrhypwidth}{\textbf{#1}\ }%
  \ifdim\jcrhypwidth<0.5in\setlength{\jcrhypwidth}{0.5in}\fi
  \leavevmode\hspace*{-0.5in}\makebox[\jcrhypwidth][l]{\textbf{#1}}\ignorespaces}

% ---- Front matter: paragraphs without indentation ----
\newenvironment{jcrfront}{\par\begingroup\setlength{\parindent}{0pt}}{\par\endgroup}

% ---- References: half-inch hanging indent ----
\ifdefined\cslhangindent
  \setlength{\cslhangindent}{0.5in}
\else
  \newlength{\cslhangindent}
  \setlength{\cslhangindent}{0.5in}
\fi

% ---- Statistics helpers ----
\newcommand{\Mstat}[1]{\ensuremath{M = #1}}
\newcommand{\SDstat}[1]{\ensuremath{\text{SD} = #1}}
\newcommand{\pval}[1]{\ensuremath{p = #1}}
\newcommand{\Fstat}[3]{\ensuremath{F(#1, #2) = #3}}
\newcommand{\tstat}[2]{\ensuremath{t(#1) = #2}}
\newcommand{\chisq}[2]{\ensuremath{\chi^2(#1) = #2}}
\newcommand{\CI}[2]{95\% CI [#1, #2]}

% ---- MANUSCRIPT SETUP (submission format) ----
% Double spaced, left justified, no page numbers, headers, or footers. The
% line spacing comes from the "linestretch" option (2 by default, which gives
% 23 lines on a page, as Word's double spacing does in JCR's sample
% manuscripts; 1.655 is LaTeX's tighter \doublespacing).
% Primary headings: bold, centered, all caps, an extra double space before
% and after. Secondary: flush left. Tertiary: indented, italic, run in, with
% a period.
\newcommand{\jcrsetupmanuscript}{%
  \pagestyle{empty}%
  \thispagestyle{empty}%
  \raggedright%
  \setlength{\parindent}{0.5in}%
  \titleformat{\section}%
    {\normalfont\normalsize\bfseries\centering\MakeUppercase}{}{0pt}{}%
  \titlespacing*{\section}{0pt}{\baselineskip}{\baselineskip}%
  \titleformat{\subsection}%
    {\normalfont\normalsize}{}{0pt}{}%
  \titlespacing*{\subsection}{0pt}{\baselineskip}{\baselineskip}%
  \titleformat{\subsubsection}[runin]%
    {\normalfont\normalsize\itshape}{}{0pt}{}[.]%
  \titlespacing*{\subsubsection}{\parindent}{\baselineskip}{0.5em}%
  % JCR uses three levels; deeper headings are set like the third
  \titleformat{\paragraph}[runin]%
    {\normalfont\normalsize\itshape}{}{0pt}{}[.]%
  \titlespacing*{\paragraph}{\parindent}{\baselineskip}{0.5em}%
  \titleformat{\subparagraph}[runin]%
    {\normalfont\normalsize\itshape}{}{0pt}{}[.]%
  \titlespacing*{\subparagraph}{\parindent}{\baselineskip}{0.5em}%
}

% ---- ARTICLE SETUP (publication-style preview; not for submission) ----
\newcommand{\jcrsetuparticle}{%
  \pagestyle{fancy}%
  \fancyhf{}%
  \fancyfoot[C]{\small\thepage}%
  \renewcommand{\headrulewidth}{0pt}%
  \fancypagestyle{plain}{%
    \fancyhf{}%
    \fancyfoot[C]{\small\thepage}%
    \renewcommand{\headrulewidth}{0pt}%
  }%
  \thispagestyle{plain}%
  \setlength{\parindent}{1.5em}%
  \titleformat{\section}%
    {\normalfont\large\bfseries\itshape\raggedright}{}{0pt}{}%
  \titlespacing*{\section}{0pt}{18pt}{6pt}%
  \titleformat{\subsection}%
    {\normalfont\normalsize\bfseries\itshape\raggedright}{}{0pt}{}%
  \titlespacing*{\subsection}{0pt}{14pt}{4pt}%
  \titleformat{\subsubsection}%
    {\normalfont\normalsize\itshape\raggedright}{}{0pt}{}%
  \titlespacing*{\subsubsection}{0pt}{10pt}{3pt}%
  \titleformat{\paragraph}%
    {\normalfont\normalsize\itshape\raggedright}{}{0pt}{}%
  \titlespacing*{\paragraph}{0pt}{10pt}{3pt}%
  \titleformat{\subparagraph}%
    {\normalfont\normalsize\itshape\raggedright}{}{0pt}{}%
  \titlespacing*{\subparagraph}{0pt}{10pt}{3pt}%
  \captionsetup{font=small,labelfont={bf},textfont={}}%
}
